\section{Experimental Setup}
\label{sec:experimental-setup}

The primary experiment is the supervisor-provided Tesla EV graph-recovery run
stored under \path{results/tesla_ev_graph_recovery}. The thesis analyzes the
saved experimental evidence rather than introducing a new benchmark run.

\subsection{Models}

The evaluated models are local Ollama models:
\texttt{deepseek-r1:8b} and \texttt{llama3.1:8b}. The supervisor notebook
records both as 4.9 GB models and labels their families as DeepSeek and Llama,
respectively. Table~\ref{tab:graph-models} summarizes the model configuration
used in the primary result matrix.

\begin{table}[!htbp]
  \centering
  \small
  \setlength{\tabcolsep}{3pt}
  \begin{tabularx}{\textwidth}{@{}L{32mm}L{24mm}L{22mm}Y@{}}
    \toprule
    Model & Family & Size in result CSV & Role in comparison \\
    \midrule
    \texttt{deepseek-r1:8b} & DeepSeek & 4.9 GB & Reasoning-oriented open-source model evaluated for structured graph recovery. \\
    \texttt{llama3.1:8b} & Llama & 4.9 GB & Same-size open-source baseline model for comparison. \\
    \bottomrule
  \end{tabularx}
  \caption{Models evaluated in the supervisor-provided Tesla EV result matrix. The size and family fields are taken from the long-format evaluation CSV.}
  \label{tab:graph-models}
\end{table}

The notebook calls Ollama through \path{/api/generate}. It sets temperature to
0.2, top-\(p\) to 0.9, and context length to 8192 in the notebook
configuration. It also defines a per-model output-token cap for
\texttt{deepseek-r1:8b} to prevent reasoning traces from running without a
known bound. These settings are documented as part of the supervisor notebook;
the result artifacts, rather than a new run, are the numerical source of
truth.

\subsection{Dataset and Run Grid}

The active dataset list contains one dataset:
\texttt{tesla\_ev}, labeled \emph{Tesla Full EV}. The active model list
contains two models. The validated result matrix contains five strategies.
Therefore, the primary evaluation matrix has ten rows:

\begin{equation}
  1 \text{ dataset} \times 2 \text{ models} \times 5 \text{ strategies}
  = 10 \text{ model-strategy cells}.
\end{equation}

Each cell is represented by one predicted adjacency matrix in the result
directory. The JSON sidecars report \texttt{n\_runs = 1}. The folder name
contains \texttt{10x}, and the notebook includes repeated-run code, but the
validated primary files used for this thesis do not provide repeated-run
metric rows. The methodology therefore treats the experiment as a
single-output-per-cell comparison.

\begin{table}[!htbp]
  \centering
  \small
  \setlength{\tabcolsep}{3pt}
  \begin{tabularx}{\textwidth}{@{}L{45mm}Y@{}}
    \toprule
    Configuration item & Value used in primary thesis evidence \\
    \midrule
    Dataset key & \texttt{tesla\_ev} \\
    Dataset label & Tesla Full EV \\
    Use-case name & \texttt{UC\_Tesla\_Energy\_Consumption} \\
    Selected features & Six: \texttt{electricenergyremaining}, \texttt{electricremaining}, \texttt{speed\_kmh}, \texttt{acceleration\_kmh2}, \texttt{mileage}, \texttt{distance\_m} \\
    Models & \texttt{deepseek-r1:8b}, \texttt{llama3.1:8b} \\
    Strategies in primary results & Five \\
    Evaluation rows & Ten model-strategy cells \\
    Per-cell repetitions in JSON sidecars & \texttt{n\_runs = 1} \\
    Primary result table & \path{evaluation_matrix_tesla_ev_long.csv} \\
    \bottomrule
  \end{tabularx}
  \caption{Primary experimental configuration used in the thesis. The table intentionally avoids unsupported repeated-run claims.}
  \label{tab:primary-experiment-config}
\end{table}

\subsection{Saved Per-Cell Artifacts}

For each evaluated model-strategy cell, the result directory contains three
main artifacts:

\begin{itemize}
  \item a \texttt{.json} sidecar with the dataset key, model, strategy, feature
  list, parsed edges, metric scores, total latency, \texttt{n\_runs}, errors,
  and raw model text;
  \item a \texttt{.npy} file storing the predicted adjacency matrix used by the
  evaluation builder;
  \item a \texttt{.png} graph visualization of the predicted edge set.
\end{itemize}

The result directory also contains aggregate artifacts: the full evaluation
matrix, a long-format evaluation table, per-model and F1 heatmaps, a long edge
export, a JSON-LD edge export, and a consensus edge table. These artifacts
allow the results chapter to move between numerical ranking, qualitative edge
inspection, and consensus analysis without recomputing generations.

\subsection{Reproducibility Boundary}

The reproducibility standard is artifact traceability rather than experiment
rerun. Every numerical claim in the methodology and results must be traceable
to one of the following:

\begin{itemize}
  \item \path{evaluation_matrix_tesla_ev.csv};
  \item \path{evaluation_matrix_tesla_ev_long.csv};
  \item per-cell result JSON files under \path{tesla_ev/<model>/<strategy>.json};
  \item \path{Ground_Truth_UC_Tesla_Energy_Consumption_signed.json};
  \item \path{edges_kg.csv};
  \item \path{consensus_tesla_ev.csv}.
\end{itemize}

This boundary prevents the thesis from importing claims from commented
notebook cells, unexecuted repeated-run sections, or the older summarization
and governance-oriented experiment. The graph-recovery artifacts are primary;
previous repository infrastructure is supplementary.

\subsection{Known Methodological Limitations}

The experimental setup has several limitations that affect interpretation.
First, the evaluation uses one Tesla EV dataset and one six-feature graph, so
the result is not a general benchmark for all vehicle systems. Second, the
ground-truth graph is physics-informed and suitable for evaluation, but it is
not a formally identified causal model from interventions. Third, the
validated artifacts provide one output per model-strategy cell, so output
variance and self-consistency cannot be quantified from the primary evidence.
Fourth, the parser enforces a strict feature vocabulary; useful but
out-of-vocabulary prose does not become a correct edge.

These limitations do not invalidate the comparison. They define its scope:
the thesis evaluates prompt decomposition and instruction structure as
practical methods for improving LLM-assisted telematics graph recovery in a
controlled, auditable Tesla EV setting.
